\chapter{经典积分定理}
\label{chap:classical-integral-theorems}

Newton--Leibniz 公式把区间内部的导数积分化为两个端点的差。Green、Gauss
和 Stokes 定理把同一原则分别写在平面区域、空间区域和曲面上：内部微分
经过积分后，只留下带有诱导定向的边界项。边界的光滑性和定向不是附属约定，
而是公式符号的一部分。

本章先在可由有限个 $C^1$ 图形分片描述的紧区域上证明公式。更一般的分片
光滑区域可通过有限分割归入这一类；相邻分片的公共边界以相反方向出现并
相消。

\section{Green 定理}

平面有界区域 $D$ 的边界取正向：沿边界行进时，区域位于左侧。若
$D=\{(x,y):a\le x\le b,\alpha(x)\le y\le\beta(x)\}$，其中
$\alpha,\beta$ 分段 $C^1$，则由一元基本定理
\[
 \iint_D\partial_yP\,dxdy
 =\int_a^b[P(x,\beta(x))-P(x,\alpha(x))]dx
 =-\int_{\partial D}P\,dx.
\]
对 $x$-方向的简单区域同理得到 $\iint_D\partial_xQ=\int_{\partial D}Qdy$。

\begin{theorem}[Green 定理]
\label{thm:green-classical}
设 $D\subset\R^2$ 是有限个上述简单区域的并，内部互不相交，边界分段
$C^1$；$P,Q$ 在 $\overline D$ 的邻域属于 $C^1$。则
\[
 \int_{\partial D}P\,dx+Q\,dy
 =\iint_D(\partial_xQ-\partial_yP)\,dxdy.
\]
有洞区域的外边界取逆时针，内边界取顺时针。
\end{theorem}

\begin{proof}
对同时为两种简单区域的情形，将上面的两个一元基本定理相加。一般分片区域
逐片应用；每条内部公共弧在两个分片中方向相反，线积分相消，只余外边界。
若有洞，内边界也是剩余边界，但其诱导方向与外边界相反。
\end{proof}

取 $P=-y/2,Q=x/2$ 得面积公式
\[
 |D|=\frac12\int_{\partial D}(x\,dy-y\,dx).
\]
删去分片或边界正则性后，公式未必仍能在 Riemann 意义下表述；后续测度论
和弱导数理论会给出更广的版本。

\section{Gauss 散度定理}

对 $F=(F_1,F_2,F_3)$，散度为
$\operatorname{div}F=\partial_xF_1+\partial_yF_2+\partial_zF_3$。

\begin{theorem}[Gauss 散度定理]
\label{thm:gauss-divergence-classical}
设 $\Omega\subset\R^3$ 是可分成有限个坐标方向简单区域的有界区域，边界
分片 $C^1$，$n$ 为外单位法向。若 $F$ 在 $\overline\Omega$ 的邻域属于
$C^1$，则
\[
 \iiint_\Omega\operatorname{div}F\,dV
 =\iint_{\partial\Omega}F\cdot n\,dS.
\]
\end{theorem}

\begin{proof}
只证 $F=(0,0,R)$ 的分量公式。若
$\Omega=\{(x,y,z):(x,y)\in D,\alpha(x,y)\le z\le\beta(x,y)\}$，
Fubini 定理和一元基本定理给出
\[
 \iiint_\Omega\partial_zR\,dV
 =\iint_D[R(x,y,\beta)-R(x,y,\alpha)]dxdy.
\]
上、下图面的有向面积向量的第三分量分别为 $1$ 和 $-1$，侧面的第三分量
为零，故右端正是 $\iint_{\partial\Omega}F\cdot n\,dS$。另外两个分量
循环置换坐标。有限分片后，内部曲面的法向相反，通量相消。
\end{proof}

散度描述单位体积的净流出率。取 $F(x)=x/3$ 便得到
$|\Omega|=\frac13\int_{\partial\Omega}x\cdot n\,dS$。

定理要求向量场在区域的一个邻域中光滑，不能只检查边界附近。令
\[
 F(x)=\frac{x}{\norm{x}^3}\qquad(x\ne0).
\]
直接计算得 $\operatorname{div}F=0$ 于 $\R^3\setminus\{0\}$，但在半径
$R$ 的球面上，$F\cdot n=1/R^2$，故外向通量恒为 $4\pi$。若球包含原点，
散度定理的光滑性假设不成立；这个非零通量不能与区域内的零散度相矛盾。
删去半径 $\rho$ 的小球后应用定理，内球面的诱导法向指向原点，其通量为
$-4\pi$，两条边界的贡献恰好相消。

\section{经典 Stokes 定理}

对 $F=(P,Q,R)$，旋度定义为
\[
 \operatorname{curl}F=
 (R_y-Q_z,\ P_z-R_x,\ Q_x-P_y).
\]

\begin{theorem}[经典 Stokes 定理]
\label{thm:stokes-classical-surface}
设 $S=X(D)$ 是定向分片 $C^2$ 曲面，$D$ 为 Green 定理适用的平面区域，
边界定向与法向满足右手规则。若 $F$ 在 $S$ 的邻域属于 $C^1$，则
\[
 \int_{\partial S}F\cdot dx
 =\iint_S(\operatorname{curl}F)\cdot n\,dS.
\]
\end{theorem}

\begin{proof}
把边界线积分拉回参数域。令
$A(u,v)=F(X(u,v))\cdot X_u$、
$B(u,v)=F(X(u,v))\cdot X_v$。Green 定理给出
\[
 \int_{\partial D}A\,du+B\,dv
 =\iint_D(\partial_uB-\partial_vA)\,du\,dv.
\]
链式法则展开后，$F\cdot X_{uv}$ 项相消，余项恰为
\[
 \partial_uB-\partial_vA
 =\operatorname{curl}F(X)\cdot(X_u\times X_v).
\]
这就是所需曲面积分。分片曲面仍由内部边界相消得到。
\end{proof}

\section{势函数与路径无关}

\begin{theorem}[保守场判据]
\label{thm:conservative-field-path-independence}
设 $U\subset\R^n$ 开且道路连通，$F:U\to\R^n$ 连续。以下条件等价：
\begin{enumerate}[label=\textup{(\roman*)}]
 \item 存在 $C^1$ 函数 $\varphi$ 使 $F=\nabla\varphi$；
 \item 任意分段 $C^1$ 曲线的积分只依赖端点；
 \item 每条分段 $C^1$ 闭曲线上的积分为零。
\end{enumerate}
\end{theorem}

\begin{proof}
(i) 推出 (ii) 是链式法则，(ii) 推出 (iii) 显然。设 (iii) 成立，固定
$x_0\in U$，对 $x\in U$ 选道路 $\gamma_x$ 从 $x_0$ 到 $x$，定义
$\varphi(x)=\int_{\gamma_x}F\cdot dx$。任意两条道路首尾拼接成闭路，故
定义与道路无关。在 $x$ 附近取完全落在 $U$ 中的短线段，则
\[
 \varphi(x+he_j)-\varphi(x)=\int_0^hF_j(x+te_j)dt.
\]
一元基本定理给出 $\partial_j\varphi=F_j$，所以 $\varphi\in C^1$。
\end{proof}

若 $U\subset\R^2$ 单连通且 $P_y=Q_x$，先对分段 $C^1$ 简单闭曲线应用
Green 定理；单连通性保证曲线围成的有界区域仍包含于 $U$。一般闭路可先
在有限个自交点处分割成简单闭路，再由分段逼近处理。于是所有闭路积分为零。
对有洞区域结论会失败：在
$\R^2\setminus\{0\}$ 上
\[
 F(x,y)=\frac{1}{x^2+y^2}(-y,x)
\]
满足 $Q_x-P_y=0$，但沿单位圆积分为 $2\pi$，所以不存在全局势函数。
局部无旋与全局保守之间缺少的是定义域的拓扑条件。

在星形域上，可以不经过闭路分解而直接写出势函数。以下结论也是下一章
Poincaré 引理在一次形式情形的坐标版本。

\begin{proposition}[星形域上的势函数]
\label{prop:star-shaped-potential}
设 $U\subset\R^n$ 关于 $a\in U$ 星形，$F\in C^1(U;\R^n)$，并且
\[
 \partial_iF_j=\partial_jF_i\qquad(1\le i,j\le n).
\]
则
\[
 \varphi(x)=\int_0^1F(a+t(x-a))\cdot(x-a)\,dt
\]
满足 $\nabla\varphi=F$。
\end{proposition}

\begin{proof}
记 $z=x-a$。对第 $i$ 个变量求导，并把对称偏导条件用在第一项，得到
\begin{align*}
 \partial_i\varphi(x)
 &=\int_0^1\left[
 t\sum_{j=1}^n\partial_iF_j(a+tz)z_j+F_i(a+tz)\right]dt\\
 &=\int_0^1\left[
 t\sum_{j=1}^n\partial_jF_i(a+tz)z_j+F_i(a+tz)\right]dt\\
 &=\int_0^1\frac{d}{dt}\bigl[tF_i(a+tz)\bigr]dt
 =F_i(x).
\end{align*}
被积函数及其偏导在每条紧线段上连续，因而这里可以按第
\chapterref{chap:uniform-convergence-preservation}中含参数积分的初等结论
逐项求导；也可直接对差商使用一致连续性。
\end{proof}

\section*{习题}
\addcontentsline{toc}{section}{习题}
\subsection*{A组　Green 定理}
\begin{enumerate}[label=\textbf{\thechapter.\arabic*.}]
 \exerciseitem{green-simple-domain-proof}
 对同时为两种简单区域的 $D$ 完整写出 Green 定理证明。
 \exerciseitem{green-annulus-orientation}
 在圆环上验证内外边界的定向，并用 Green 定理计算一个具体线积分。
 \exerciseitem{green-area-formula}
 推导三种面积公式 $\int x\,dy$、$-\int y\,dx$ 和正文的对称形式。
 \exerciseitem{green-internal-boundary-cancellation}
 证明有限分片时公共边界积分相消，并说明三片交于一点不产生额外项。
\end{enumerate}
\subsection*{B组　Gauss 与 Stokes}
\begin{enumerate}[label=\textbf{\thechapter.\arabic*.},resume]
 \exerciseitem{gauss-box-direct}
 对长方体直接证明散度定理，并与六个面的通量逐项比较。
 \exerciseitem{gauss-volume-formula}
 用 $F=x/3$ 推导体积公式，并计算球和圆锥的体积。
 \exerciseitem{gauss-inverse-square-flux}
 计算 $F(x)=x/\norm{x}^3$ 穿过包围原点的球面的通量，并解释为何不能直接在
 含原点区域上应用散度定理。
 \exerciseitem{stokes-parameter-expansion}
 展开 $\partial_u(F(X)\cdot X_v)-\partial_v(F(X)\cdot X_u)$，验证旋度恒等式。
 \exerciseitem{stokes-hemisphere-disk}
 分别用上半球和其边界圆盘计算同一闭曲线积分，说明 Stokes 定理只依赖边界。
 \exerciseitem{curl-gradient-divcurl}
 对 $C^2$ 函数和向量场证明 $\operatorname{curl}\nabla\varphi=0$ 及
 $\operatorname{div}\operatorname{curl}F=0$。
\end{enumerate}
\subsection*{C组　势函数与条件}
\begin{enumerate}[label=\textbf{\thechapter.\arabic*.},resume]
 \exerciseitem{path-independent-potential-construction}
 补全保守场判据中势函数连续可微性的证明。
 \exerciseitem{punctured-plane-circulation}
 核对穿孔平面向量场处处无旋，并计算其单位圆环流。
 \exerciseitem{star-shaped-curl-free-potential}
 在星形域中直接定义
 $\varphi(x)=\int_0^1F(tx)\cdot x\,dt$，证明 $C^1$ 无旋场满足
 $\nabla\varphi=F$。
 \optionalexerciseitem{classical-theorem-hypothesis-audit}
 比较三条经典积分定理对区域、边界和被积场正则性的要求，列出不能由本章
 证明覆盖的边界例。
\end{enumerate}
